% Documentclass Settings
	\documentclass[a4paper, twoside]{article}
	% Margins
	\usepackage{geometry}
		\geometry{
			% showframe,
			left   = 20mm,
			right  = 20mm,
			top    = 20mm,
			bottom = 20mm
		}

	% Headers and footers
	\usepackage{fancyhdr}
		\pagestyle{fancy}
		\fancyhf{}
		\fancyhead[LE, RO]{\thepage}
		\fancyhead[CE]{\textsc{Cong Wen}}
		\fancyhead[CO]{\textsc\rightmark}

		\setlength{\headheight}{15pt}
		\renewcommand\headrulewidth{0pt}

	\usepackage[dvipsnames]{xcolor}    % Color
	\usepackage[colorlinks, citecolor = green, linkcolor = blue]{hyperref}





\input{../universal-pre}

\title{\tbf{Kollar1986}}
\author{\tbf{Notes} by Cong Wen, 10/2022}
\date{}
% Last Updated: July 11, 2022



\begin{document}
\maketitle
\thispagestyle{empty}
\begin{abstract}
	This is my self-use notes\footnote{Last updated on October 04, 2022} on \tbf{Higher direct images of dualizing sheaves}\cite{Kollar1986-1,Kollar1986-2}. The original papers have been very well written so this note really is only for personal use. Please reach out if you find anything incorrect.
\end{abstract}
\tableofcontents
% ----------------------------------------------------------------------------------------------------


\begin{mtv}
	To compute sheaf cohomology on a smooth variety $X$, it is natural to consider a surjective map of varieties $f:X\ar Y$ and reduce the computation to $Y$ via Leray spectral sequence. A starting point would be the properties of $R^{j}f\zS\cO_{X}$ and $R^{j}f\zS\oga_{X}$.
\end{mtv}

\begin{thm}
	Let $X$ be a smooth projective variety, $Y$ any projective variety, $f:X\ar Y$ surjective, then
	\begin{enr}
		\item $R^{j}f\zS\oga_{X}, \fal j\gqs0$ is torsion-free,
		\item if $L$ is ample on $Y$, then
		\[
			\mathrm{H}_{}^{i}(Y,L\ot R^{j}f\zS\oga_{X})=0, \fal i>0.
		\]
	\end{enr}
	Moreover, if $Y^{0}\sbs Y$ is the smooth locus, $Y-Y^{0}$ is SNC, $f^{0}:=f|_{f\xI(Y^{0})}$, then
	\begin{align*}
		R^{j}f\zS\oga_{X/Y}&\cong{}^{u}\cF^{b}(R^{n-k+j}f^{0}\zS\uL{\bC})\\
	  	R^{j}f\zS\cO_{X}&\cong {}^{l}\Gr^{0}(R^{j}f^{0}\zS\uL{\bC})\\
		R\xB f\zS\oga_{X}&\cong R^{j}f\zS\oga_{X}[-j].
	\end{align*}

\end{thm}




% ----------------------------------------------------------------------------------------------------
\bibliographystyle{alpha}
\bibliography{../universal-ref}
\end{document}

